\begin{chapterpage}{The Kernel}
    \begin{chapteritems}
        \item Kernel API
        \item Device Table
        \item File Table 
    \end{chapteritems}
\end{chapterpage}

\section{Introduction}

Most of the work in the system is done by device drivers. The kernel is
just a thin layer that provides a common interface to the drivers.

The core device drivers required to start the system are embedded into the kernel ROM bank (bank 0).
Other device drivers, including alternatives to the core drivers, can be stored
in ROM banks 1-255 or even be loaded into RAM at runtime.

You are free to decide what device drivers are used for different tasks. If you, for example, want 
a USB keyboard instead of the default PS/2 keyboard, it is possible to load and use a custom USB 
keyboard device driver instead of the default PS/2 keyboard driver.

The kernel is inspired by Unix, where almost everything is a file. The device driver interface is, 
consequently, based on functions such as open, read, write, and close.


\section{API Functions}

The kernel's internal code and data structures are subject to change. 
Only the functions listed in this section are part of the stable interface.

\begin{longtblr}{colspec = {X[5,l] X[3,l] X[8,l] },
width = \linewidth,
rowhead = 1,
rowfoot = 0}
Name & Address & Description \\
\hline
reset               & \$c000    & System reset \\
alloc               & \$c003    & Allocates one RAM bank (8 kB) \\
free                & \$c006    & Frees one RAM bank \\          
device\_register    & \$c009    & Registers a new device driver \\
device\_unregister  & \$c00c    & Unregisters a device \\          
file\_open          & \$c00f    & Opens a file \\                   
file\_read\_byte    & \$c012    & Reads one byte from a file \\
file\_read\_block   & \$c015    & Reads a block of bytes from a file \\
file\_write\_byte   & \$c018    & Writes one byte to a file \\
file\_write\_block  & \$c01b    & Writes a block of bytes to a file \\
file\_seek          & \$c01e    & Seeks a position in a file \\
file\_position      & \$c021    & Returns current position in a file \\
file\_close         & \$c024    & Closes a file \\
file\_ioctl         & \$c027    & Driver-specific control commands \\
stdin\_read         & \$c02a    & Reads one byte from stdin \\     
stdout\_write       & \$c02d    & Writes one byte to stdout \\              
get\_ticks          & \$c030    & Returns the jiffy clock \\              
set\_ticks          & \$c033    & Sets the jiffy clock \\                
get\_errmsg         & \$c036    & Returns system error message \\
set\_errmsg         & \$c039    & Sets system error message \\
get\_info           & \$c03c    & Get system info \\
\end{longtblr}


\section{Passing and Returning Values}

Values are passed to the kernel functions and returned through 
the A, X, Y registers, and ten zero page registers 
named arg1-arg10 (address \$02-\$0b).

Multi-byte values are passed or returned least significant byte
first (little-endian).

arg1-arg10 are reserved for kernel calls and must not be used for
other purposes.


\section{Error Handling}

All kernel functions must indicate an error
condition by setting the carry flag on return.

An error code is always returned in the A register. The A register
can return other values if no error occurred. A plain text error
message is available through the get\_errmsg
function.

List of error codes:

\begin{longtblr}{colspec = {X[1,l] X[3,l] X[5,l] },
width = \linewidth,
rowhead = 1,
rowfoot = 0}
Code & Name & Description \\
\hline
0     & EOK                 & No error                                \\
1     & ENOMEM              & Out of memory                           \\
2     & EMEMINUSE           & Requested memory already in use         \\
3     & ENODEVSPC           & Too many devices                        \\
4     & ENODEV              & Device does not exist                   \\
5     & EDEVBUSY            & Device is in use                        \\
6     & ENOFILESPC          & Too many files                          \\
7     & ENOTFOUND           & File not found                          \\
8     & ENOTFILE            & Path is not a file                      \\
9     & ENOTDIR             & Path is not a directory                 \\
10    & ENOTOPEN            & File not open                           \\
11    & EFILEOPEN           & File is open                            \\
12    & EINVAL              & Invalid argument or operation           \\
13    & EIO                 & Generic driver failure                  \\
14    & EEOF                & End of file                             \\
255   & ECUSTOM             & Custom error message                    \\
\end{longtblr}


\section{Device Table}

The kernel uses a device table to keep track of device drivers. The device table has room for 16 devices, 
and contains the following fields:

\begin{longtblr}{colspec = {X[3,l] X[2,l] X[11,l] },
width = \linewidth,
rowhead = 1,
rowfoot = 0}
Field & Values & Description \\
\hline
addr\_lo    & 0-255 & Pointer to device driver code, bits 0-7 \\
addr\_hi    & 0-255 & Pointer to device driver code, bits 8-15 \\
bank        & 0-255 & The RAM or ROM bank where the driver is stored \\
flags       & 0-255 & See list of flags below \\
\end{longtblr}

The flags stored in the device table are:

\begin{itemize}
    \item Bit 0: Set if reading is supported.
    \item Bit 1: Set if writing is supported.
    \item Bit 2: Set if appending is supported.
    \item Bit 3: Set if seeking is supported.
    \item Bit 4: Set if a block device, unset if a char device.
    \item Bits 5-6: Driver location, where 00 = Not present, 01 = ROM bank 0 or low RAM,
    10 = ROM bank 1-255, 11 = RAM bank 0-255.
    \item Bit 7: Set if kernel is requested to call the driver's on\_irq function on every IRQ.
\end{itemize}

A device driver is added to the device table when you call the device\_register function. The
device driver is removed when you call device\_unregister. It is not recommended or
safe to manipulate the device table directly.


\section{File Table}

The kernel has a file table to keep track of all opened files. There is room for 16 files in
the file table. The table holds the following information:

\begin{longtblr}{colspec = {X[3,l] X[2,l] X[10,l] },
width = \linewidth,
rowhead = 1,
rowfoot = 0}
Field & Values & Description \\
\hline
driver\_id  & 0-15 & Driver's index in device table \\
context\_id & 0-255 & The context ID is used internally by drivers to keep track of multiple open files \\
flags       & 0-255 & See list of flags below \\
\end{longtblr}

The flags in the file table are:

\begin{itemize}
    \item Bit 0: Set if opened for reading.
    \item Bit 1: Set if opened for writing.
    \item Bit 2: Set if opened for appending.
    \item Bit 3: Set if a file, unset if a directory.
    \item Bits 4-6: Reserved.
    \item Bit 7: Set if file is open.
\end{itemize}

A file is added to the file table when you call the file\_open function, and it is removed when
you call the file\_close function. It is not recommended or safe to manipulate the file table directly.

Some file handles have a fixed meaning to the system:

\begin{itemize}
    \item 0: STDIN.
    \item 1: STDOUT.
    \item 2: Keyboard.
    \item 3: Mouse.
    \item 4: SNES game controllers.
\end{itemize}


\section{API Function Specifications}

\subsection{Function: reset}

\textbf{Call address:} \$c000 \\
\textbf{Input:} None \\
\textbf{Output:} None \\
\textbf{Errors:} None

This function does a cold reset of the system. It is called automatically when
the system is reset. It is also possible to call the function while
running the system.

The function closes any open files, unregisters any current device drivers, and
frees all allocated RAM banks.

Finally, it registers all built-in base drivers that are needed to use
the system:

\begin{enumerate}\addtocounter{enumi}{-1}
    \item X16 terminal driver.
    \item Generic I2C driver.
    \item PS/2 keyboard driver.
    \item PS/2 mouse driver.
    \item SNES game controller driver.
    \item SD card driver.
\end{enumerate}


\subsection{Function: alloc}

\textbf{Call address:} \$c003 \\
\textbf{Input:} A = RAM bank \\
\textbf{Output:} A = RAM bank \\
\textbf{Errors:} ENOMEM, EMEMINUSE

The X16 has no hardware memory management. The kernel's memory management functions 
(alloc and free) keep track of which part of banked RAM is free to use. All well-behaved
applications should use these functions. They may otherwise overwrite custom
device drivers loaded into banked RAM.

RAM bank 0 is used by the kernel and can't be allocated for other purposes.

If this function is called with A = 0, it will search for and allocate the next available
RAM bank. Error ENOMEM is set if there is no free memory.

If the function is called with A > 0, it will attempt to allocate the specified RAM bank.
EMEMINUSE is returned if that bank is already allocated.

A successfully allocated RAM bank is returned in the A register.


\subsection{Function: free}

\textbf{Call address:} \$c006 \\
\textbf{Input:} A = RAM bank \\
\textbf{Output:} None \\
\textbf{Errors:} EINVAL

This function frees a RAM bank that was allocated with the alloc function.

RAM bank 0 is used by the system and may not be freed. If this function is
called with A = 0, the EINVAL error is returned.


\subsection{Function: device\_register}

\textbf{Call address:} \$c009 \\
\textbf{Input:} Y = Device ID, arg1-2 = Driver address, arg3 = Driver bank \\
\textbf{Output:} A = Device ID \\
\textbf{Errors:} ENODEVSPC, EDEVBUSY, EIO

The function registers a device driver that is already present in memory.

The system's base drivers needed to boot the system are embedded into
the kernel ROM bank, and are automatically registered and initialized by
the reset function.

If called with Y = 255, the first free device ID will be used. Otherwise, the function 
attempts to use the specified device ID.

The driver's on\_register function is automatically called when calling this function. 
The kernel's device table is updated with information on the new driver, 
including flags returned by the on\_register function.

The ENODEVSPC error is returned if the driver table is full.

The EDEVBUSY error is returned if the specified device ID is already in use.
In that case you must first call device\_unregister to free the current driver.

EIO is a driver-specific error.


\subsection{Function: device\_unregister}

\textbf{Call address:} \$c00c \\
\textbf{Input:} Y = Device ID \\
\textbf{Output:} None \\
\textbf{Errors:} EDEVBUSY, EIO

Unregisters the specified device driver.

The EDEVBUSY error is returned if the device driver is linked to an open file.

EIO is a driver-specific error.


\subsection{Function: file\_open}

\textbf{Call address:} \$c00f\\
\textbf{Input:} X = File handle, Y = Device ID, A = Flags, arg1-2 = Path address, arg3 = Path bank \\
\textbf{Output:} X = File handle, Y = Context ID \\
\textbf{Errors:} ENODEV, ENOFILESPC, EFILEOPEN, ENOTFILE, ENOTFOUND, EINVAL, EIO

Opens a file using the specified device driver.

If the X register is set to \$ff, the new file is linked to the next available file handle. 
Otherwise, the function attempts to use the specified file handle.

The file is opened using the device driver specified in the Y register.

The following flags are supported:

\begin{itemize}   
    \item Bit 0: Set if opened for reading.
    \item Bit 1: Set if opened for writing.
    \item Bit 2: Set if opened for appending.
\end{itemize}

You must set exactly one flag. The kernel doesn't support opening a file
in multiple modes (or no mode).

The path must be a null-terminated string.

If successfully opened, the file is entered into the file table. The file
handle is returned in the X register and the driver's file context ID is
returned in the Y register. The context is managed by the driver, and is used 
by the driver to keep track of multiple files that are opened at the same time.

ENODEV is returned if no driver is linked to the specified device ID.

The ENOFILESPC error is set if the file table is full. There can be at most 16 open files
at any given time. The same error is also returned if the number of files the driver
can handle is exceeded.

The EFILEOPEN error is returned if you try to link a new file to a file handle that is already in use.

ENOTFILE is returned if attempting to write or append to a directory.

The function fails with ENOTFOUND if the file is not found.

The EINVAL error is returned if the device driver doesn't have the capability
requested in the flags argument. The same error is returned if not exactly one flag is set.

EIO is a driver-specific error.


\subsection{Function: file\_read\_byte}

\textbf{Call address:} \$c012\\
\textbf{Input:} X = File handle \\
\textbf{Output:} A = Value \\
\textbf{Errors:} ENOTOPEN, EINVAL, EEOF, EIO

Reads one byte from the specified file.

The ENOTOPEN error is returned if the specified file handle is not open.

The EINVAL error is set if the file is not opened for reading.

EEOF is returned if the read operation was attempted at end of the file and 
no byte value was returned. A successful read operation of the final byte does 
not return EEOF; the caller receives EEOF on the following read attempt.

EIO is a driver-specific error.


\subsection{Function: file\_read\_block}

\textbf{Call address:} \$c015\\
\textbf{Input:} X = File handle, arg1-2 = Requested number of bytes, arg3-4 = Buffer address, arg5 = Buffer bank \\
\textbf{Output:} arg6-7 = Number of bytes read \\
\textbf{Errors:} ENOTOPEN, EINVAL, EEOF, EIO

Reads a block from the specified file, at most the requested number of bytes. It is not guaranteed that
you get the requested number of bytes.

The ENOTOPEN error is returned if the specified file handle is not open.

The EINVAL error is returned if the file is not opened for reading.

EEOF is set if the read operation was attempted at end of the file and 
no byte value was returned. A successful read operation of the final byte does 
not return EEOF; the caller receives EEOF on the next read attempt.

EIO is a driver-specific error.


\subsection{Function: file\_write\_byte}

\textbf{Call address:} \$c018\\
\textbf{Input:} X = File handle, A = Byte value \\
\textbf{Output:} None \\
\textbf{Errors:} ENOTOPEN, ENOTFILE, EINVAL, EIO

Writes one byte to the specified file.

The ENOTOPEN error is returned if the specified file handle is not open.

ENOTFILE is returned if attempting to write or append
to a directory.

The EINVAL error is returned if the file is not opened for writing or appending.

EIO is a driver-specific error.


\subsection{Function: file\_write\_block}

\textbf{Call address:} \$c01b\\
\textbf{Input:} X = File handle, arg1-2 = Buffer size, arg3-4 = Buffer address, arg5 = Buffer bank \\
\textbf{Output:} arg6-7 = Number of bytes written \\
\textbf{Errors:} ENOTOPEN, ENOTFILE, EINVAL, EIO

Writes one block to the specified file, at most the entire buffer size passed in arg1-2.

The ENOTOPEN error is returned if the specified file handle is not open.

ENOTFILE is returned if attempting to write or append
to a directory.

The EINVAL error is returned if the file is not opened for writing or appending.

EIO is a driver-specific error.


\subsection{Function: file\_seek}

\textbf{Call address:} \$c01e\\
\textbf{Input:} X = File handle, arg1-4 = File position \\
\textbf{Output:} arg1-4 = File position \\
\textbf{Errors:} ENOTOPEN, EINVAL, EIO

Moves the file pointer to the requested position.

The file position after this call is returned in arg1-4.

The kernel only allows writing to files in a sequential order at the end of the file, 
and seeking is thus allowed only on files opened for reading. Attempting to 
seek on a file opened for writing or appending returns 
the EINVAL error.

EINVAL is also returned if the device driver doesn't support
seeking.

The ENOTOPEN error is returned if the specified file handle is not open.

EIO is a driver-specific error.


\subsection{Function: file\_position}

\textbf{Call address:} \$c021\\
\textbf{Input:} X = File handle \\
\textbf{Output:} arg1-4 = File position \\
\textbf{Errors:} ENOTOPEN, EINVAL, EIO

Returns the current position in the specified file in arg1-4.

ENOTOPEN is returned if the file is not open.

EINVAL can be returned if the driver doesn't support
reading the file positions.

EIO is a driver-specific error.


\subsection{Function: file\_close}

\textbf{Call address:} \$c024\\
\textbf{Input:} X = File handle \\
\textbf{Output:} None \\
\textbf{Errors:} ENOTOPEN, EIO

Closes the specified file and removes it from the file table.

ENOTOPEN is returned if the file handle is not open.

EIO is a driver-specific error.


\subsection{Function: file\_ioctl}

\textbf{Call address:} \$c027\\
\textbf{Input:} X = File handle, A = Device type, arg1 = Command, arg2 = Input length, arg3-4 = Input data or input buffer address, arg5 = Input buffer bank,
arg6 = Output max length, arg7-8 = Output buffer address, arg9 = Output buffer bank \\
\textbf{Output:} A = Output length, arg7-8 = Output data or output buffer address, arg9 = Output buffer bank \\
\textbf{Errors:} ENOTOPEN, EINVAL, EIO

This function executes a device-specific control command.

The file handle is given in order to let the command affect a specific file context. 
Device drivers may, however, ignore the file context if it's not relevant to the command.
ENOTOPEN is returned if the file is not open. 

If this function is called with X = \$ff, the driver must ignore it and not
return the ENOTOPEN error.

The type of device that the command is meant for is specified in the A register. EINVAL 
is returned if the driver doesn't support commands for the specified device type.

The control command is specified in arg1. The two most significant bits are used to flag if the input 
(arg3-4) and output (arg7-8) are direct values or buffer pointers. Direct values can be at most two bytes long.

\begin{itemize}
    \item Bit 7: Set if output is a direct value, clear if output is a buffer pointer.
    \item Bit 6: Set if input is a direct value, clear if input is a buffer pointer.
\end{itemize}

The input is specified in arg1-5, and the output is specified in arg7-8.

The function returns the output length in the A register. If the output is a direct value, it is returned
in arg7-8. Otherwise, the output buffer pointer and bank are returned in arg7-9.

EIO is a driver-specific error.


\subsection{Function: stdin\_read}

\textbf{Call address:} \$c02a\\
\textbf{Input:} None \\
\textbf{Output:} A = Byte value \\
\textbf{Errors:} ENOTOPEN, EINVAL, EEOF, EIO

Reads one byte from STDIN, i.e. file handle 0.

The ENOTOPEN error is returned if STDIN is not open.

The EINVAL error is set if STDIN is not opened for reading.

EEOF is returned if the read operation was attempted at end of the file and 
no byte value was returned. A successful read operation of the final byte does 
not return EEOF; the caller receives EEOF on the following read attempt. This
doesn't normally happen when STDIN is connected to a keyboard that
has an endless stream conceptually.

EIO is a driver-specific error.


\subsection{Function: stdout\_write}

\textbf{Call address:} \$c02d\\
\textbf{Input:} A = Byte value \\
\textbf{Output:} None \\
\textbf{Errors:} ENOTOPEN, EINVAL, EIO

Writes one byte to STDOUT, i.e. file handle 1.

The ENOTOPEN error is returned if STDOUT is not open.

The EINVAL error is returned if STDOUT is not opened for writing or appending.

EIO is a driver-specific error.


\subsection{Function: get\_ticks}

\textbf{Call address:} \$c030\\
\textbf{Input:} None \\
\textbf{Output:} X, Y, A = Counter value \\
\textbf{Errors:} None

Returns the jiffy clock tick counter that is
incremented on every IRQ (60 times per second).


\subsection{Function: set\_ticks}

\textbf{Call address:} \$c033\\
\textbf{Input:} X, Y, A = Counter value \\
\textbf{Output:} None \\
\textbf{Errors:} None

Sets the jiffy clock tick counter that is
incremented on every IRQ (60 times per second).


\subsection{Function: get\_errmsg}

\textbf{Call address:} \$c036\\
\textbf{Input:} None \\
\textbf{Output:} A = Error number, arg1-2 = Error message address, arg3 = Error message bank \\
\textbf{Errors:} None

Returns the last error number in the A register and a pointer to
the last error message, that is always a null-terminated string.


\subsection{Function: set\_errmsg}

\textbf{Call address:} \$c039\\
\textbf{Input:} A = Error number, arg1-2 = Error message address, arg3 = Error message bank \\
\textbf{Output:} None \\
\textbf{Errors:} None

Sets the error number and message that can be retrieved with the
get\_errmsg function.


\subsection{Function: get\_info}

\textbf{Call address:} \$c03c\\
\textbf{Input:} None \\
\textbf{Output:} arg1-3 = Kernel version (major, minor patch), arg4-5 = String address, arg6 = String bank \\
\textbf{Errors:} None

Returns the kernel version, and a plain null-terminated text description of the system.
